\sectionguard
\section{The Historicity Dispute as a Live Question}

\subsection{Why this section does not adjudicate}

The question whether four apparitions occurred at Tepeyac in December
1531 and whether a Nahua man called Juan Diego Cuauhtlatoatzin
received them is disputed among competent, credentialed, believing
scholars, and it is disputed on grounds that are historical rather
than theological. Nothing in the Church's acts settles it as a
historical question; nothing in the surviving record settles it
either way beyond reasonable disagreement; and no study produced by
an AI contributor without archival access has any business
pretending otherwise.

This section therefore does three things. It states each position as
its strongest advocates state it. It identifies the actual points on
which the disagreement turns, so that a reader can see that the two
sides are not simply talking past one another. And it records, with
precision, what this study has and has not verified about each
position, because the scholars involved are living or recently
deceased professionals whose views deserve to be reported at the
level at which they were actually checked, and not one word beyond.

\subsection{The two positions, stated at their strongest}

\subsubsection*{The case for a late origin}

The strongest form of the sceptical case is not that anyone lied. It
is that a devotion of genuine sixteenth-century antiquity acquired,
in the seventeenth century, an origin narrative of a familiar and
well-attested literary kind, and that the narrative was received with
such rapidity and such joy that its recent composition was forgotten
within two generations.

On this account the sequence runs: a hermitage of Our Lady is
established at Tepeyac, a place of pre-conquest cult; by the
mid-1550s an image in that hermitage attracts a rapidly growing
Spanish and then Indian devotion, ignited by reports of a cure; the
Franciscans object publicly in 1556 that the image is recent, human
work and that the miracle reports are unverified; the archbishop
promotes the cult regardless; the shrine prospers for a century
without any narrative of an apparition; in 1648 a Mexican-born
diocesan priest supplies one, in a Creole-patriotic theological key;
in 1649 the chaplain of the shrine supplies a Nahuatl version for
Indian readers; and from 1666 onward the tradition is defended,
juridically documented, and finally made the object of Roman
liturgical concession.

The load-bearing evidence is the silence of \S3.1, the 1556 file of
\S3.2, the viceroy's letter of 1575, and the absence of any dated
sixteenth-century manuscript of the narrative. The load-bearing
argument is that the coincidence of all four is not the kind of thing
that accident produces.

\subsubsection*{The case for the tradition}

The strongest form of the apparitionist case is likewise not what its
opponents sometimes describe. It is not an appeal to piety against
evidence. It is a claim about \emph{which} evidence a historian of
sixteenth-century Mexico is entitled to privilege.

On this account the sceptical case rests almost entirely on Spanish
alphabetic sources produced by and for the colonial administration
and the Spanish clergy---exactly the sources least likely to record
an event whose significance was Indigenous, whose language was
Nahuatl, and whose principal effect was among a conquered population
in the decade immediately following a catastrophic military defeat
and in the middle of a Franciscan campaign against Indian image cult.
The proper sources are the Nahua ones: the annals, the songs, the
testaments, the pictorial records, the memory transmitted in a
culture with sophisticated non-alphabetic and oral means of
preservation. Those sources exist, they are numerous, they are early,
and---the apparitionist argues---the sceptical tradition has
systematically undervalued them because it reads them through a
European documentary standard they were never meant to satisfy.

On this account the narrative in the 1649 tract is not a
seventeenth-century invention but a sixteenth-century Nahuatl
composition that Laso de la Vega printed; the shrine's own published
dossier of Indigenous, mestizo and Spanish documents is the case in
chief (\S2.3); the 1556 controversy is evidence of a bitter
inter-jurisdictional quarrel about image cult, not of ignorance of an
apparition; and the 1666 depositions preserve, at the last possible
moment, a chain of oral transmission at Cuauhtitl\'an that a
documentary historian is not entitled to dismiss merely because it is
oral.

\subsection{The named positions and what was actually checked}

The following inventory names the principal published statements. The
right-hand column states the single argument each is chiefly known
for. Except where the note says otherwise, this study verified each
entry's \emph{bibliographic identity} through accessible catalogue
records and states the position at a reported ceiling; it did not
read the books.\endnote{Bibliographic identity of the Poole, Sousa--Poole--Lockhart,
Brading, Burkhart, von Wobeser, Gonz\'alez Fern\'andez and Ch\'avez
entries verified 2026-07-25 against Open Library catalogue records
(\url{https://openlibrary.org/search.json}) and, for the 1998 edition,
against the Internet Archive record
\url{https://archive.org/details/storyofguadalupe84lass} (Stanford
University Press and UCLA Latin American Center Publications, 1998;
UCLA Latin American Studies vol.~84 / Nahuatl Studies Series no.~5;
ISBNs 0804734828 and 0804734836). The Poole 2005 article citation
(\work{The Americas} 62:1, July 2005, pp.~1--16) was verified against
the Project MUSE and Cambridge Core records; the article text is
behind a paywall and was not read. \emph{None of the books in this
table was read for this study}; the summary in the right-hand column
states each author's position at the level at which it is reported in
the accessible literature and in the sources cited elsewhere in this
section, and no wording is attributed to any of these authors.
Miguel Le\'on-Portilla's \work{Tonantzin Guadalupe} (2000), which is
frequently cited in this debate, could not be verified in any
catalogue reachable for this study and is therefore not used.}

\begin{positionstable}
Joaqu\'in Garc\'ia Icazbalceta (d.~1894) & \work{Carta acerca del
origen de la imagen de Nuestra Se\~nora de Guadalupe de M\'exico}
(1883; printed 1896) & The silence of Zum\'arraga and of every
sixteenth-century writer who had occasion to mention the event
\emph{(read in full for this study)}\\
\'Angel Mar\'ia Garibay K. & Address to the 1960
Mariological Congress, Mexico City, printed 1961 & The
\work{Nican mopohua} is a mid-sixteenth-century Nahuatl composition
by Sahag\'un's Nahua collaborators from older notes
\emph{(read in O'Gorman's detailed exposition)}\\
Edmundo O'Gorman & \work{Destierro de sombras} (UNAM,
1986) & The cult originates in 1555--1556 under Mont\'ufar; the
\work{Nican mopohua} is Antonio Valeriano's own literary invention
for Indian readers, not a record of events
\emph{(read at length for this study)}\\
Stafford Poole, C.M. & \work{Our Lady of Guadalupe:
The Origins and Sources of a Mexican National Symbol, 1531--1797}
(Univ.\ of Arizona Press, 1995; 2nd ed.\ 2017); \work{The Guadalupan
Controversies in Mexico} (Stanford, 2006); ``History versus Juan
Diego,'' \work{The Americas} 62:1 (2005) & The apparition account is
unknown before 1648; Juan Diego's historicity is not established by
the surviving evidence\\
Lisa Sousa, Stafford Poole, James Lockhart & \work{The Story of
Guadalupe: Luis Laso de la Vega's} \work{Huei tlamahui\c{c}oltica}
\work{of 1649} (Stanford / UCLA, 1998) & Philological dependence of
the Nahuatl narrative on S\'anchez 1648, hence a seventeenth-century
composition\\
D.~A.~Brading & \work{Mexican Phoenix: Our Lady of Guadalupe. Image
and Tradition Across Five Centuries} (Cambridge, 2001) & The
tradition's theological and political formation traced as a history
of the tradition itself\\
Louise M.~Burkhart & \work{Before Guadalupe: The Virgin Mary in Early
Colonial Nahuatl Literature} (Albany, 2001) & The Nahuatl Marian
literature that precedes and surrounds the Guadalupe text\\
Gisela von Wobeser & \work{Or\'igenes del culto a Nuestra Se\~nora de
Guadalupe, 1521--1688} (UNAM-IIH / FCE, 2020) & Recent full-length
scholarly treatment of the cult's origins\\
Xavier Escalada, S.J. & \work{Enciclopedia Guadalupana}
and the publication of the \term{C\'odice Escalada} (1995--1997) &
A parchment dated 1548 bearing a Sahag\'un signature and a Valeriano
glyph as a sixteenth-century witness to Juan Diego\\
Fidel Gonz\'alez Fern\'andez, Eduardo Ch\'avez S\'anchez, Jos\'e Luis
Guerrero Rosado & \work{El encuentro de la Virgen de Guadalupe y Juan
Diego} (Porr\'ua, 1999) & The convergent Indigenous, mestizo and
Spanish documentary corpus establishes Juan Diego's historicity\\
Eduardo Ch\'avez & \work{Our Lady of Guadalupe and Saint Juan Diego:
The Historical Evidence} (Rowman \& Littlefield, 2006) & The same
case restated for an English readership by the cause's postulator\\
\end{positionstable}

Two facts about this table deserve saying aloud, because the dispute
is often narrated as a conflict between faith and scholarship, and it
is not. Garc\'ia Icazbalceta wrote his letter at his archbishop's
command and begged him to keep it private; Stafford Poole was a
Vincentian priest; O'Gorman's book was published by the national
university's historical institute. On the other side, Fidel
Gonz\'alez Fern\'andez was a consultor of the Congregation for the
Causes of Saints and Eduardo Ch\'avez was the postulator of Juan
Diego's cause and a trained historian. The line does not run between
believers and unbelievers. It runs between two defensible readings of
the same difficult corpus.

\subsection{Where the disagreement actually lives}

Four cruxes carry almost the whole weight.

\subsubsection*{Crux 1: the date of the \work{Nican mopohua}}

Everything turns on this, and it is a philological question. If the
Nahuatl narrative is a sixteenth-century composition, the case for a
seventeenth-century invention collapses, whatever remains to be said
about the events themselves. If it depends on S\'anchez's 1648
Spanish book, the case for antiquity loses its principal support.

The remarkable feature of the modern literature is that the
philological question does not track the theological one. Garibay,
who dated the text to Sahag\'un's Tlatelolco circle between 1564 and
1570, was arguing \emph{for} the tradition. O'Gorman, who accepts
that Antonio Valeriano wrote it, uses that conclusion against the
tradition: on his reading Valeriano composed it about 1556 precisely
as a work of persuasion addressed to Indians, and its literary
quality is evidence of authorship, not of testimony. The
Sousa--Poole--Lockhart edition argues for dependence on S\'anchez and
a date after 1648. A reader who wants to know whether the text is
early must therefore follow an argument about Nahuatl style,
vocabulary, and textual parallels---not about
piety.\endnote{Garibay's thesis and O'Gorman's rebuttal are set out in
O'Gorman, \work{Destierro de sombras}, Appendix~I, ``La relaci\'on de
las apariciones (\work{Nican mopohua}), supuesta obra de
colaboradores ind\'igenas de Sahag\'un,'' read 2026-07-25 at
\url{https://historicas.unam.mx/publicaciones/publicadigital/libros/222c/222c_04_05_01_ApendicePrimero.pdf};
O'Gorman there cites \'Angel Mar\'ia Garibay K., ``La maternidad
espiritual de Mar\'ia en el mensaje guadalupano,'' address of 10
October 1960 to the Mariological Congress at Mexico City, printed in
\work{La maternidad espiritual de Mar\'ia} (Mexico: Jus, 1961),
pp.~187--202, at pp.~191--194. O'Gorman's own attribution of the
\work{Nican mopohua} to Valeriano is stated in part~I, ch.~3, ``La
invenci\'on del guadalupanismo ind\'igena,'' to which his Appendix~I
refers back, and is reflected in his part~II, ch.~3, \S III
(``nuestro an\'alisis de la obra de Valeriano''). The
Sousa--Poole--Lockhart position is reported here at the level of the
edition's identity and its well-known thesis; the edition was not
read (\S4.3, note).}

\subsubsection*{Crux 2: the \term{C\'odice Escalada}}

In 1995 the Jesuit Xavier Escalada announced, and in 1997 published,
a small piece of treated animal skin bearing a drawing of an apparition
scene and a kneeling Indian, glosses in Nahuatl including the name
Juan Diego and a death date, the year 1548, what is presented as the
signature of Bernardino de Sahag\'un, and a glyph read as that of
Antonio Valeriano. If genuine it is a sixteenth-century witness to
Juan Diego's existence and death, and it would be decisive.

Its authenticity is disputed. The reported grounds of dispute include
the character of the drawing and of the Valeriano glyph, the syntax
of the glosses, the presence of two inks, and---most consequentially
for method---the conditions under which it was examined: the reported
tests of 1997 are said to have been non-destructive and partly
conducted on photographic reproductions at the owner's request, and
one of the participating scientists is reported to have publicly
qualified the certainty of the findings in 2002. This study reached
no primary publication of any examination and reports the whole
matter at that
ceiling.\endnote{Recorded negative finding. No primary publication of
the 1997 examinations, no report by any named examining institution,
and no peer-reviewed authentication or refutation of the
\term{C\'odice Escalada} was reached through the searches performed on
2026-07-25. The account above is reported from general reference and
press summaries and is stated as report only; no reader should treat
it as this study's finding about the object. The
Bas\'ilica's own \work{Documentos Hist\'oricos} page lists the item as
``C\'odice 1548 o `Escalada','' describing a codex showing an Indian
kneeling before the image
(\url{https://virgendeguadalupe.org.mx/documentos-historicos/}, read
2026-07-25); the 1998 statement of the Historical Commission
described it, as reported, as ``a direct testimony of Juan Diego's
historicity'' (\S5.3). Both are positions of interested parties and
are cited as such.}

\subsubsection*{Crux 3: what counts as a source}

The two schools weight source-types differently in a way that is
methodological rather than evidential. One holds that a narrative
first attested in print in 1648 must be treated as a 1648 text until
an earlier witness is produced, and that oral transmission recorded
in 1666 by an interested tribunal cannot bear the weight of a
1531 event. The other holds that a documentary standard calibrated to
European archives systematically erases Indigenous memory, and that
requiring a dated alphabetic Nahuatl manuscript from the 1530s is
requiring evidence that the conditions of the period made almost
impossible.

Both propositions are true of something. The dispute is about which
is true of this case.

\subsubsection*{Crux 4: the argument from silence}

Its force is proportional to how confident one is that a given writer
\emph{would} have mentioned the event. That confidence is highest for
Zum\'arraga, who is made the narrative's second protagonist, and
lowest for writers with no occasion to touch the matter. The
sceptical case is strong exactly to the degree that it concentrates
on Zum\'arraga and on the 1556 file, and weaker as it lengthens into
a list. The apparitionist case is strong exactly to the degree that it
supplies positive early evidence, and weaker as it relies on
explaining silences away.

\subsection{What a Catholic reader is entitled to conclude}

Three things, and they are enough.

First, that the shrine at Tepeyac and its Marian cult are genuinely
of the sixteenth century, on evidence that nobody disputes. The
devotion is old, Indigenous in its reception from very early, and
continuous.

Second, that the narrative of the four apparitions is transmitted to
us by witnesses of 1648--1649 and is defended by testimony of 1666,
and that whether it rests on earlier material is a genuinely open
historical question on which serious Catholic scholars differ. A
reader who finds the sceptical case persuasive is in the company of a
Vincentian priest and of the archbishop's own chosen historian; a
reader who finds the traditional case persuasive is in the company of
the postulator of the cause and of a consultor of the Congregation
for the Causes of Saints.

Third---and this is developed in \S9---that the Church's own acts do
not require the reader to settle the question, because none of them
is a determination of the historical question. What they determine is
something else, and \S1 and \S5 set out exactly what.
